\documentclass{article}
\usepackage[UTF8]{ctex}
\usepackage{amsmath,amssymb,amsfonts}
\usepackage{enumitem}
\usepackage{geometry}
\geometry{a4paper, margin=1in}
% === 错题本排版优化（不参与题目标识） ===
\setlength{\parindent}{0pt}
\setlength{\parskip}{0.35em plus 0.10em minus 0.08em}
\setlength{\emergencystretch}{3em}
\setlist{itemsep=0.25em, topsep=0.40em, parsep=0pt}
\allowdisplaybreaks[2]
% === 错题本排版优化结束 ===
\begin{document}
\section*{中国科学院大学2026年高等代数真题}
\section*{15．中国科学院大学2026年高等代数试卷（回忆版）}
（注．题目来自考生回忆，部分表述或数据可能有误）

\subsection{题目1}
1．（15 分）设 \(f(x), g(x), h(x)\) 为实系数多项式，试给出＂存在实系数多项式 \(u(x), v(x)\) 使得\(u(x) f(x)+v(x) g(x)=h(x)^{n}\) 的充要条件，并求出满足条件下的 \(u(x), v(x)\) 的表达式。

\subsection{题目2}
2．（20 分）设 \(A=\left(a_{i j}\right)_{n \times n}\) 是 \(n\) 阶实对称正定矩阵，其特征值 \(t_{1}<t_{2}<\cdots<t_{n}\) ，考虑下面二次型

\[
f\left(x_{1}, x_{2}, \cdots, x_{n}\right)=\left|\begin{array}{ccccc}
a_{11} & a_{12} & \cdots & a_{1 n} & x_{1} \\
a_{21} & a_{22} & \cdots & a_{2 n} & x_{2} \\
\vdots & \vdots & \vdots & & \vdots \\
a_{n 1} & a_{n 2} & \cdots & a_{n n} & x_{n} \\
x_{1} & x_{2} & \cdots & x_{n} & a\left(x_{1}^{2}+x_{2}^{2}+\cdots+x_{n}^{2}\right)
\end{array}\right| .
\]

试确定 \(a\) 的取值，使二次型 \(f\left(x_{1}, x_{2}, \cdots, x_{n}\right)\) 分别正定，负定，不定，写出解与秩。

\subsection{题目3}
3．（15 分）设 \(A, B\) 为 \(n\) 阶正交矩阵，且 \(\operatorname{det}(A)+\operatorname{det}(B)=0\) ，证明：\(A+B\) 非可逆．

\subsection{题目4}
4．（25 分）设 \(\mathcal{A}\) 为 \(n\) 维线性空间 \(V\) 上的一个幂等线性变换，即满足 \(\mathcal{A}^{2}=\mathcal{A}\) 。证明：

\begin{itemize}
  \item[（1）] \(\operatorname{Ker} \mathcal{A}=\{X-\mathcal{A} X \mid X \in V\}\) ．
  \item[（2）] \(V=\operatorname{Ker} \mathcal{A} \oplus \operatorname{Im} \mathcal{A}\) ．
  \item[（3）] 设 \(\mathcal{B}\) 是 \(V\) 上的另一个线性变换，那么 \(\operatorname{Ker} \mathcal{A}, \operatorname{Im} \mathcal{A}\) 都是 \(\mathcal{B}\) 的不变子空间当且仅当 \(\mathcal{A} \mathcal{B}=\mathcal{B} \mathcal{A}\) ．
\end{itemize}

\subsection{题目5}
5．（20 分）设 \(\mathcal{A}, \mathcal{B}\) 为 \(n\) 维复向量空间 \(V\) 上的线性变换，且满足 \(\mathcal{A} \mathcal{B}-\mathcal{B} \mathcal{A}=2026 \mathcal{B}\) ，证明：存在 \(V\)的一组基，使得 \(\mathcal{A}, \mathcal{B}\) 在这组基下的矩阵均为下三角矩阵。

\subsection{题目6}
6．（20 分）设 \(V_{1}, V_{2}, \cdots, V_{m}\) 为数域 \(\mathbb{F}\) 上的 \(n\) 维向量空间 \(V\) 的任意 \(m\) 个 \(r\) 维子空间，证明：存在 \(V\)的子空间 \(W\) ，使得

\[
V=V_{1} \oplus W=V_{2} \oplus W=\cdots=V_{m} \oplus W
\]

\subsection{题目7}
7．（15 分）设 \(M \in M_{n}(\mathbb{C}), \bar{M}\) 为 \(M\) 的复共轭矩阵。证明：若 \(M\) 与 \(\bar{M}\) 相似，则 \(M\) 复相似于一个实矩阵。

\subsection{题目8}
8．（20 分）设（V，\(((-,-))\) 是 \(n\) 维欧氏空间，\(\left(e_{1}, e_{2}, \cdots, e_{n}\right)\) 和 \(\left(e^{1}, e^{2}, \cdots, e^{n}\right)\) 为 \(V\) 的两组基，满足

\[
\left(e_{i}, e^{j}\right)=\delta_{i j}(i, j=1,2, \cdots, n) .
\]

证明：若 \(\left(e_{i}, e_{j}\right) \leq 0\) 对所有的 \(i \neq j\) ，则 \(\left(e^{i}, e^{j}\right) \geq 0\) 对所有的 \(i \neq j\) ．
\end{document}
